\emph{(i)} $x$ is a scalar and commutes with $\rho$, and $\rho^{2} = 0$ (Proposition \claimref{std-algebra.radical-ideal-and-quotient}: the radical generates a nilpotent ideal, here with $J^{2} = 0$), so $(x + \rho)^{k} = x^{k} + kx^{k-1}\rho$ and $u(x + \rho) = u(x) + u'(x)\rho$ for a polynomial $u$. For $f = u/h$ with $h(x) \ne 0$, $h(x + \rho) = h(x) + h'(x)\rho$ has the inverse $h(x)^{-1} - h'(x)h(x)^{-2}\rho$, and multiplying gives $f(x) + (u'h - uh')(x)h(x)^{-2}\rho = f(x) + f'(x)\rho$. The check runs $f = (x^{3} + 3x)/(x^{2} + 1)$, with $f' = (x^{4} + 3)/(x^{2} + 1)^{2}$, at four rational points against hand-stated values.

\emph{(ii)} The $b_{i}$ commute, so the multinomial expansion of $\varepsilon^{k}$ holds, and every monomial with a repeated factor vanishes since $b_{i}^{2} = 0$; the surviving monomials are the products over $k$-element subsets, each appearing $k!$ times, so $\varepsilon^{k} = k!\,e_{k}(b)$, and $e_{N+1} = 0$. Expanding $\prod(1 + b_{i})$ gives $\sum_{k}e_{k}(b)$, which is $\sum_{k \le N}\varepsilon^{k}/k!$, and the exponential series of the nilpotent $\varepsilon$ stops at $k = N$. For the realisations: $\rho_{2i-1}\rho_{2i}$ is even, so it commutes with every product of other generators, and $(\rho_{2i-1}\rho_{2i})^{2} = -\rho_{2i-1}^{2}\rho_{2i}^{2} = 0$; for mutually orthogonal null vectors, $n_{j}n_{l} = -n_{l}n_{j}$ for $j \ne l$ and $n_{j}^{2} = 0$ (Proposition \claimref{std-algebra.witt-basis-car}), so the same two computations apply. In $\Cl(2N,2N)$ the vectors $n_{j} = \tfrac12(e_{j} + e_{2N+j})$, $j = 1, \dots, 2N$, are such a family, a maximal isotropic set. The check verifies the identities for $N = 2, \dots, 6$ in $\Cl(0,0,2N)$ and for $N = 3$ with six orthogonal null vectors of $\Cl(6,6)$.

\emph{(iii)} $Bl_{+} = e_{t}e_{r}e_{t} + e_{t}e_{r}e_{r} = -e_{r} - e_{t} = -l_{+}$ and $l_{+}B = e_{t}e_{t}e_{r} + e_{r}e_{t}e_{r} = e_{r} + e_{t} = l_{+}$, so $B$ anticommutes with $l_{+}$ and $l_{+}\widetilde{R} = l_{+}(\alpha - \beta B) = (\alpha + \beta B)l_{+} = Rl_{+}$; hence $Rl_{+}\widetilde{R} = R^{2}l_{+} = (\alpha^{2} + \beta^{2} + 2\alpha\beta B)l_{+} = (\alpha - \beta)^{2}l_{+}$, and likewise $Rl_{-}\widetilde{R} = (\alpha + \beta)^{2}l_{-}$ since $Bl_{-} = l_{-}$. For central $d$ with $d^{2} = 0$, $\exp(\tfrac12 dB) = 1 + \tfrac12 dB$ by Proposition \claimref{std-algebra.bivector-exponential} (the series stops), which is $\alpha = 1$, $\beta = \tfrac12 d$, and $(1 \mp \tfrac12 d)^{2} = 1 \mp d$. In $\Cl(1,1,2)$, $d = s\rho_{1}\rho_{2}$ is even and built from generators orthogonal to $e_{t}, e_{r}$, so it is central in the subalgebra generated by $e_{t}, e_{r}, \rho_{1}, \rho_{2}$ and squares to $0$. For the finite boost only $R^{2} = \cosh\phi + \sinh\phi\,B$ enters, and $\cosh\phi - \sinh\phi = \tfrac54 - \tfrac34 = \tfrac12$, $\cosh\phi + \sinh\phi = 2$; $(\tfrac54)^{2} - (\tfrac34)^{2} = 1$.

\emph{(iv)} A rotor of the plane is $R = c + sI$ with $c^{2} + s^{2} = 1$, so $R^{2} = 2cR - 1$ and $(1 + R)^{2} = 1 + 2R + 2cR - 1 = 2(1 + c)R$; and $(1 + R)(1 + R)^{\sim} = 2 + R + \widetilde{R} = 2(1 + c)$, positive for $R \ne -1$. So $S = (1 + R)/\lvert 1 + R\rvert$ is a rotor with $S^{2} = R$, and its scalar part satisfies $c'^{2} = (1 + c)/2$ while $2c's' = s$. From $c_{0} = 0$, $s_{0} = 1$: $c_{1}^{2} = \tfrac12$, so $c_{1} = \tfrac12 r_{1}$ with $r_{1}^{2} = 2$, and inductively $c_{k+1}^{2} = (1 + \tfrac12 r_{k})/2 = \tfrac14(2 + r_{k})$, so $c_{k+1} = \tfrac12 r_{k+1}$ with $r_{k+1}^{2} = 2 + r_{k}$ (positive roots throughout). For Archimedes' recursion write $c, s$ for $c_{k}, s_{k}$ and $c', s'$ for $c_{k+1}, s_{k+1}$, with $s = 2c's'$ and $c = 2c'^{2} - 1$: then $U_{k} + L_{k} = 2^{k+1}s(1 + c)/c$ and $2U_{k}L_{k}/(U_{k} + L_{k}) = 2^{k+2}s/(1 + c) = 2^{k+2}\cdot 2c's'/(2c'^{2}) = 2^{k+2}s'/c' = U_{k+1}$; and $U_{k+1}L_{k} = 2^{k+2}(s'/c')\,2^{k+1}\cdot 2c's' = (2^{k+2}s')^{2} = L_{k+1}^{2}$. The check implements the field $\Q(r_{1}, r_{2}, \dots)$ exactly, as iterated quadratic extensions with $r_{k+1}^{2} = 2 + r_{k}$, and verifies $R_{k+1}^{2} = R_{k}$, $R_{k}\widetilde{R_{k}} = 1$ for $k = 0, \dots, 5$ in $\Cl(2,0)$, the identity $(1 + R)^{2} = 2(1 + c)R$ at four rational rotors, and the recursion for $k = 1, \dots, 5$. With the exponential chart of Proposition \claimref{std-algebra.bivector-exponential}, $R_{0} = I = \exp(\tfrac{\pi}{2}I)$ and $R_{k} = \exp(\theta_{k}I)$ with $\theta_{k} = \pi/2^{k+1}$ (the square root with positive scalar part), so $c_{k} = \cos\theta_{k}$, $s_{k} = \sin\theta_{k}$, and $\sin\theta < \theta < \tan\theta$ on $(0, \tfrac{\pi}{2})$ gives $L_{k} < \pi < U_{k}$; $\tan\theta - \sin\theta = \tfrac12\theta^{3} + O(\theta^{5})$ gives $U_{k} - L_{k} = 2^{k+1}(\tfrac12\theta_{k}^{3} + O(\theta_{k}^{5}))$, so $4^{k}(U_{k} - L_{k}) \to \pi^{3}/8$. Numerically $(L_{6}, U_{6}) = (3.14128, 3.14222)$ and $4^{6}(U_{6} - L_{6}) = 3.876$ against $\pi^{3}/8 = 3.8758$. These analytic statements are not part of the check.

\emph{(v)} Computed exactly: $\tan = \sin\cdot\cos^{-1}$ and $\arctan = \int(1 + x^{2})^{-1}$ as truncated power series over $\Q$, through degree $21$; every denominator is odd (the largest is $194896477400625 = 3^{9}\cdot 5^{4}\cdot 7^{3}\cdot 11\cdot 13\cdot 17\cdot 19$, and the coefficients of $\arctan$ are $\pm 1/(2k + 1)$), while $p \mid k!$ for every prime $p \le k$.
